\documentclass[12pt,a4paper]{article}

\usepackage{amsmath,amssymb,amsthm}
\usepackage{hyperref}
\usepackage{natbib}
\usepackage{geometry}
\usepackage{booktabs}
\usepackage{xcolor}
\usepackage{microtype}
\usepackage{mdframed}

\geometry{margin=1in}

\hypersetup{
  colorlinks=true,
  linkcolor=blue,
  citecolor=blue,
  urlcolor=blue,
  pdftitle={The Geometric Origin of the Bekenstein-Hawking Entropy Coefficient}
}

\newtheorem{proposition}{Proposition}
\newtheorem{corollary}{Corollary}

\title{\textbf{The Geometric Origin of the Bekenstein--Hawking\\
Entropy Coefficient}\\[0.5em]
\large Why $\eta = 1/(4\ell_P^2)$ Is Not a Free Parameter
in Jacobson's Thermodynamic Derivation of Gravity}

\author{Heath W.\ Mahaffey\\
\small Independent Researcher, Entiat, Washington 98822, USA\\
\small \href{mailto:hmahaffeyges@gmail.com}{hmahaffeyges@gmail.com}\\
\small \url{https://github.com/hmahaffeyges/IAM-Validation}}

\date{April 2026\\
\small Zenodo DOI:
\href{https://doi.org/10.5281/zenodo.18702042}{10.5281/zenodo.18702042}}

\begin{document}
\maketitle

\begin{abstract}
Jacobson's 1995 derivation of the Einstein field equations from
thermodynamics makes two assumptions: entropy proportional to horizon
area, $S = \eta A$, and that $\eta$ is a free parameter identified
with the observed Bekenstein--Hawking value $1/(4\ell_P^2)$ by matching
to Newton's constant. We show that neither assumption is independent
of the causal geometry of flat spacetime.

The proportionality $S \propto A$ follows necessarily from the
irreversibility of quantum-to-classical transitions (decoherence): each
such transition produces classical information that must be encoded at
the nearest causally accessible surface. Causal boundaries are
two-dimensional --- volume encoding is causally forbidden because
information inside the horizon is inaccessible to external observers
--- so entropy scales with area, not volume. The assumption $S \propto A$
is therefore a causal necessity, not an independent postulate.

The coefficient $\eta = 1/(4\ell_P^2)$ is geometrically determined
within Jacobson's framework by the ratio of two quantities derivable
from flat spacetime independently of $S \propto A$: the Einstein
equation normalization $8\pi$ (from Gauss's law in three dimensions
giving $4\pi$, doubled by the Bianchi identity trace structure) and
the Euclidean Rindler cone angle $2\pi$ (the unique topological period
eliminating the conical singularity at the horizon). Their ratio:
\begin{equation*}
\frac{1}{4} = \frac{2\pi}{8\pi}
\end{equation*}
fixes $\eta$ as a geometric consequence. The same ratio determines
the minimum horizon area per decoherence event to be $4\ell_P^2$,
consistent with exactly one bit per $4\ell_P^2$ --- three statements
of the same geometry. Jacobson's two assumptions reduce to one
derived result. The single remaining open question is why a minimum
resolvable length $\ell_P = \sqrt{\hbar G/c^3}$ exists at all ---
a quantum gravity microstructure question that thermodynamic
approaches cannot address from within their own framework.
\end{abstract}

\noindent\textbf{Keywords:} Bekenstein--Hawking entropy; Jacobson
thermodynamic gravity; Rindler horizon; Unruh effect; holographic principle;
IAM's Law; geometric origin of Newton's constant

% ─────────────────────────────────────────────────────────────────────────────
\section{The Problem}
\label{sec:problem}
% ─────────────────────────────────────────────────────────────────────────────

Jacobson~\citep{Jacobson1995} showed that the Einstein field equations
follow from the Clausius relation $\delta Q = T\,dS$ applied to local
Rindler horizons. The derivation requires two inputs:

\begin{enumerate}
\item The Unruh temperature of the Rindler horizon~\citep{Unruh1976,
Davies1975}:
\begin{equation}
T = \frac{\hbar\kappa}{2\pi k_B},
\label{eq:unruh}
\end{equation}
where $\kappa$ is the surface gravity.

\item Entropy proportional to horizon area:
\begin{equation}
S = \eta A,
\label{eq:S_eta_A}
\end{equation}
where $\eta$ is treated as a free parameter.
\end{enumerate}

Applying the Clausius relation and requiring it to hold for all null
vectors yields the Einstein equations with Newton's constant:
\begin{equation}
G = \frac{c^4}{4\hbar\eta}.
\label{eq:G_eta}
\end{equation}
Jacobson then identifies $\eta = c^3/(4\hbar G) = 1/(4\ell_P^2)$ by
matching Eq.~\eqref{eq:G_eta} to the observed value of $G$. The
coefficient $1/4$ is recovered, not derived.

This gap has been shared by every thermodynamic approach to gravity
since~\citep{CaiKim2005,Padmanabhan2010}. Loop quantum gravity recovers
$\eta = 1/(4\ell_P^2)$ only by fixing the Barbero--Immirzi
parameter~\citep{Rovelli1995}. String theory recovers it in specific
limits~\citep{Strominger1996}. Neither approach derives it from the
causal geometry of the horizon alone.

The present paper shows that both assumptions are geometric
consequences of the causal structure of flat spacetime. Neither
is independent. Section~\ref{sec:area} derives $S \propto A$
from the irreversibility of quantum-to-classical transitions.
Section~\ref{sec:factors} derives the coefficient $1/4$ from
the ratio $2\pi/8\pi$.

% ─────────────────────────────────────────────────────────────────────────────
\section{Why Entropy Scales with Area: The Causal Accessibility Argument}
\label{sec:area}
% ─────────────────────────────────────────────────────────────────────────────

Jacobson assumed $S \propto A$ without derivation. We show it follows
from three physical facts that hold independently of any entropy-area
postulate.

\subsection{Decoherence Produces Classical Information}

Every irreversible quantum-to-classical transition --- every instance
of a quantum system evolving from superposition to a definite classical
outcome through interaction with its environment --- produces classical
information. This is the decoherence process~\citep{Zurek2003}. A
quantum system $\mathcal{S}$ interacting with environment $\mathcal{E}$
evolves as:
\begin{equation}
|\psi\rangle_\mathcal{S} \otimes |E_0\rangle
\;\longrightarrow\;
\sum_i c_i |s_i\rangle \otimes |E_i\rangle,
\label{eq:decoherence}
\end{equation}
where $\{|s_i\rangle\}$ are the pointer states selected by the
interaction. When the environmental states $\{|E_i\rangle\}$ become
orthogonal, the reduced density matrix of $\mathcal{S}$ becomes
diagonal in the pointer basis:
\begin{equation}
\rho_\mathcal{S} = \mathrm{Tr}_\mathcal{E}|\Psi\rangle\langle\Psi|
\;\longrightarrow\; \sum_i |c_i|^2 |s_i\rangle\langle s_i|.
\label{eq:diagonal}
\end{equation}
The system has selected a definite classical outcome. This
transition is irreversible: recovering the off-diagonal elements
would require erasing the produced information from the environment,
which by Landauer's principle~\citep{Landauer1961} dissipates
$k_B T \ln 2$ per bit into the environment and increases its entropy.
Each decoherence event produces $I = \log_2 N$ bits of classical
information, where $N$ is the number of distinguishable pointer
states.

\subsection{Classical Information Must Be Encoded at a Causally Accessible Surface}

Classical information produced by a decoherence event must be
permanently encoded somewhere accessible to observers outside the
system. The constraint is causal: information is accessible to an
external observer if and only if it lies on or outside the causal
boundary of the system. Inside the horizon, no information can reach
external observers by definition.

The causal boundary of a region is a two-dimensional surface ---
the horizon. The region itself is three-dimensional. Therefore:

\begin{itemize}
\item \textbf{Surface (area) encoding:} information on the
2D causal boundary is accessible to all external observers.
\item \textbf{Volume encoding:} information inside the 3D bulk
is causally inaccessible to external observers.
\end{itemize}

If the produced information is to be permanently accessible --- as
it must be, since decoherence is irreversible and the information
cannot be destroyed without violating the second law --- it must be
encoded at the 2D causal boundary, not in the 3D bulk.

\subsection{Area Scaling Is a Causal Necessity}

Since every irreversible quantum-to-classical transition encodes its
produced information on the 2D causal boundary, and since the
information content of the boundary grows with the number of
independent encoding sites on that surface, the total information
encoded on the boundary scales with the boundary's area. Therefore:
\begin{equation}
S \propto A.
\label{eq:S_propto_A_derived}
\end{equation}
This is not an independent postulate. It follows from:
(1) decoherence produces classical information irreversibly
(Eq.~\ref{eq:decoherence}--\ref{eq:diagonal}),
(2) permanent information must be encoded at the causally accessible
boundary, and
(3) the boundary is a 2D surface whose information capacity scales
with area.

Jacobson's assumption $S = \eta A$ is therefore not an additional
input beyond the causal geometry of spacetime. It is a consequence
of the irreversibility of quantum-to-classical transitions and the
causal structure of the horizon. The only remaining question is
the value of $\eta$ --- the proportionality constant --- which
is derived in the following section.

% ─────────────────────────────────────────────────────────────────────────────
\section{The Two Geometric Factors Determining $\eta$}
\label{sec:factors}
% ─────────────────────────────────────────────────────────────────────────────

The coefficient $4$ in $G = c^4/(4\hbar\eta)$ is the ratio:
\begin{equation}
4 = \frac{8\pi}{2\pi}.
\label{eq:four}
\end{equation}
Both $8\pi$ and $2\pi$ are geometric facts about flat spacetime. Neither
assumes entropy proportional to area.

\subsection{The $2\pi$: Euclidean Rindler Cone Angle}

The Rindler metric near an accelerating horizon takes the form:
\begin{equation}
ds^2 = -\kappa^2\rho^2\,dt^2 + d\rho^2 + d\mathbf{x}_\perp^2.
\label{eq:rindler}
\end{equation}
Under Euclidean continuation $t \to -i\tau$:
\begin{equation}
ds^2_E = \kappa^2\rho^2\,d\tau^2 + d\rho^2 + d\mathbf{x}_\perp^2.
\label{eq:euclidean_rindler}
\end{equation}
In the $(\rho, \kappa\tau)$ plane, this is a cone with $\rho$ as the
radial coordinate and $\kappa\tau$ as the angular coordinate. At $\rho = 0$
(the horizon), this cone has a conical singularity unless $\kappa\tau$ has
period exactly $2\pi$. This is the unique condition for a smooth, flat
cone with no conical deficit:
\begin{equation}
\kappa\tau \sim \kappa\tau + 2\pi
\quad \Longrightarrow \quad
\tau \sim \tau + \frac{2\pi}{\kappa}.
\label{eq:period}
\end{equation}
This is the KMS (Kubo--Martin--Schwinger) period of the Rindler
Green's function~\citep{Unruh1976,Bisognano1976}. The Unruh temperature
$T = \hbar\kappa/(2\pi k_B)$ follows directly: the thermal period
$\beta = 1/(k_B T) = 2\pi/(\hbar\kappa)$ matches the geometric period
of the Euclidean cone.

\textbf{The $2\pi$ is a topological necessity.} It is the full angle
of the Euclidean horizon cone, fixed by the requirement that the horizon
be a non-singular surface in flat spacetime. It makes no reference to
$\eta$, to $S \propto A$, or to Newton's constant.

\subsection{The $8\pi$: Einstein Equation Normalization}

The Einstein field equations are:
\begin{equation}
G_{ab} + \Lambda g_{ab} = \frac{8\pi G}{c^4}\,T_{ab}.
\label{eq:einstein}
\end{equation}
The coefficient $8\pi$ in Eq.~\eqref{eq:einstein} is fixed by two
independent geometric requirements:

\textbf{(a) Gauss's law in three spatial dimensions.}
In the Newtonian limit, the Einstein equations must reduce to Poisson's
equation:
\begin{equation}
\nabla^2\Phi = 4\pi G\rho.
\label{eq:poisson}
\end{equation}
The factor $4\pi$ is the solid angle of the unit sphere in
$\mathbb{R}^3$:
\begin{equation}
\Omega_3 = \oint d\Omega = 4\pi.
\label{eq:solid_angle}
\end{equation}
This is a purely geometric fact about three-dimensional space.

\textbf{(b) The Bianchi identity trace factor.}
The Einstein tensor $G_{ab} = R_{ab} - \frac{1}{2}Rg_{ab}$ differs
from the Ricci tensor by the trace term $\frac{1}{2}Rg_{ab}$. In the
weak-field limit, this trace structure introduces a factor of 2 relative
to the naive identification with Poisson's equation~\citep{Wald1984}.

Combining: $8\pi = 2 \times 4\pi$. Both factors are purely geometric.
The $8\pi$ makes no reference to $\eta$, to $S \propto A$, or to the
numerical value of $G$.

% ─────────────────────────────────────────────────────────────────────────────
\section{The Derivation}
\label{sec:derivation}
% ─────────────────────────────────────────────────────────────────────────────

\begin{mdframed}
\begin{proposition}[Geometric determination of the entropy coefficient]
\label{prop:main}
The entropy-area proportionality constant $\eta$ in Jacobson's
thermodynamic derivation of gravity is not a free parameter. It is
geometrically determined by:
\begin{equation}
\eta = \frac{1}{4\ell_P^2},
\qquad
\frac{1}{4} = \frac{2\pi}{8\pi},
\label{eq:eta_geometric}
\end{equation}
where $2\pi$ is the Euclidean Rindler cone angle and $8\pi$ is the
Einstein equation normalization, both derivable from the causal
structure of flat spacetime without assuming $S \propto A$.
\end{proposition}
\end{mdframed}

\medskip

\noindent\textit{Derivation.}
Jacobson's derivation yields the relation:
\begin{equation}
\frac{\hbar\eta}{2\pi} = \frac{c^4}{8\pi G}.
\label{eq:jacobson_core}
\end{equation}
This equation follows from applying $\delta Q = T\,\delta S$ to the
Rindler horizon and requiring the result to hold for all null vectors.
The $2\pi$ on the left enters through the Unruh temperature
(Eq.~\ref{eq:unruh}). The $8\pi$ on the right is the Einstein
normalization (Eq.~\ref{eq:einstein}).

Solving for $\eta$:
\begin{equation}
\eta = \frac{c^4}{8\pi G} \cdot \frac{2\pi}{\hbar}
= \frac{c^4}{4\hbar G}
= \frac{c^3}{4\hbar G / c}
= \frac{1}{4} \cdot \frac{c^3}{\hbar G}
= \frac{1}{4\ell_P^2}.
\label{eq:eta_derived}
\end{equation}
The coefficient $1/4 = 2\pi/8\pi$ is the ratio of the Euclidean
horizon cone angle to the Einstein normalization. Both are geometric.
Neither assumes $S \propto A$.

The entropy-area relation $S = \eta A$ with $\eta = 1/(4\ell_P^2)$
is therefore not a free assumption in Jacobson's framework --- it is
the unique relation consistent with the causal geometry of flat
spacetime.
\hfill$\square$

\medskip

\noindent\textbf{A note on what this shows.}
This is not a prediction of $\eta$ from outside Jacobson's framework,
nor a claim that $\eta$ could have been derived without Jacobson's
machinery. It is a demonstration that $\eta$ is geometrically fixed
\textit{within} that framework by the same causal geometry that
determines $T$ and $G_{ab}$. Jacobson had all the ingredients: the
$2\pi$ from the Euclidean horizon geometry and the $8\pi$ from the
Einstein normalization. Their ratio $2\pi/8\pi = 1/4$ fixes $\eta$
without any observational input. The matching to the observed $G$ is
then a consistency check, not an independent determination of $\eta$.

% ─────────────────────────────────────────────────────────────────────────────
\section{Independence of the Two Factors}
\label{sec:independence}
% ─────────────────────────────────────────────────────────────────────────────

The derivation requires that neither factor assumes the conclusion.

\textbf{The $2\pi$ from the Euclidean Rindler cone angle} is a
topological property of flat Minkowski spacetime near an accelerating
horizon. The Euclidean continuation of the Rindler metric produces a
cone in the $(\rho, \tau)$ plane. Regularity at the tip requires the
angular period to be $2\pi$ --- exactly the period that eliminates the
conical deficit. This is a statement about flat spacetime geometry,
derivable without reference to black holes, entropy, or $S \propto A$.
The Bisognano--Wichmann theorem~\citep{Bisognano1976} establishes this
in the full algebraic QFT setting.

\textbf{The $8\pi$ from the Einstein normalization} follows from
requiring general relativity to reduce to Newtonian gravity in the
weak-field, non-relativistic limit. The factor $4\pi$ is Gauss's
law in three spatial dimensions; the factor 2 is the trace structure
of the Einstein tensor. Both are geometric properties of three-space
and four-dimensional Riemannian geometry respectively. Neither
references entropy, Planck's constant, or $S \propto A$.

\textbf{Independence from each other:} The $2\pi$ is derived from
quantum field theory in flat spacetime via the KMS condition. The
$8\pi$ is derived from the differential geometry of the weak-field
Einstein equations. They arise from completely different physical
contexts and were identified by different historical routes --- Unruh
(1976) and Einstein (1915) respectively. Their ratio appearing as the
coefficient of $\eta$ is a structural consequence of applying
thermodynamics to a spacetime with both a causal horizon and a
Newtonian limit.

\textbf{Uniqueness:} The ratio $2\pi/8\pi = 1/4$ is not tunable.
Both factors are fixed by geometric necessities of flat spacetime.
The value $1/4$ is selected by the geometry, not chosen to match
observation.

% ─────────────────────────────────────────────────────────────────────────────
\section{Connection to Known Results}
\label{sec:connections}
% ─────────────────────────────────────────────────────────────────────────────

\subsection{The Minimum Encoding Area: One Bit = $4\ell_P^2$}

The same geometric ratio $8\pi/2\pi = 4$ that determines $\eta$
also determines the minimum horizon area disturbed by one
decoherence event. This provides an independent confirmation that
the three results --- $S \propto A$, $\eta = 1/(4\ell_P^2)$, and
the minimum encoding area --- are three expressions of the same
causal geometry.

The Landauer cost of one decoherence event (one bit of classical
information produced) at the local horizon temperature $T_H$ is:
\begin{equation}
E_\mathrm{bit} = k_B T_H = \frac{\hbar\kappa}{2\pi},
\label{eq:landauer_bit}
\end{equation}
where we use $k_B T_H = \hbar\kappa/(2\pi)$ from the Unruh
temperature (Eq.~\ref{eq:unruh}) in natural units.

This energy, deposited on the horizon, disturbs a minimum area
$\delta A_\mathrm{min}$ through the gravitational coupling. The
area response of a horizon to an energy perturbation is set by
the Einstein equation normalization:
\begin{equation}
\delta A_\mathrm{min} = \frac{G}{c^4} \cdot E_\mathrm{bit}
\cdot 8\pi,
\label{eq:area_response}
\end{equation}
where the factor $8\pi$ is the Einstein normalization
(Eq.~\ref{eq:einstein}) converting energy to area perturbation.

Substituting Eq.~\eqref{eq:landauer_bit} into
Eq.~\eqref{eq:area_response}:
\begin{equation}
\delta A_\mathrm{min}
= \frac{G}{c^4} \cdot \frac{\hbar\kappa}{2\pi} \cdot 8\pi
= \frac{4G\hbar\kappa}{c^4}.
\label{eq:dA_step}
\end{equation}
At the minimum surface gravity $\kappa_\mathrm{min} = c^2/\ell_P$
(the Planck-scale horizon, where quantum and gravitational effects
are simultaneously relevant):
\begin{equation}
\delta A_\mathrm{min}
= \frac{4G\hbar}{c^4} \cdot \frac{c^2}{\ell_P}
= \frac{4G\hbar}{c^2\ell_P}
= 4\ell_P^2,
\label{eq:dA_min}
\end{equation}
where we used $\ell_P^2 = \hbar G/c^3$ in the last step.

The minimum horizon area disturbed by one decoherence event is
$4\ell_P^2$. The entropy density $\eta = 1/(4\ell_P^2)$ gives
exactly one bit per $4\ell_P^2$:
\begin{equation}
\eta \cdot \delta A_\mathrm{min}
= \frac{1}{4\ell_P^2} \cdot 4\ell_P^2 = 1
\quad \text{(bit per encoding area).}
\label{eq:one_bit}
\end{equation}
The factor of 4 in the minimum encoding area arises from
$8\pi/(2\pi) = 4$ --- the same ratio that determines $\eta$.
This is not a numerical coincidence. It is the same geometric
statement expressed twice: the ratio of the Einstein normalization
to the Euclidean horizon period determines both the entropy
coefficient and the minimum encoding area, and they are
consistent by construction.

\subsection{Jacobson's Machine Fully Specified}

Eq.~\eqref{eq:jacobson_core} shows that Jacobson's derivation has the
structure:
\begin{equation}
\underbrace{\frac{\hbar\eta}{2\pi}}_{\text{Rindler geometry}}
= \underbrace{\frac{c^4}{8\pi G}}_{\text{Einstein geometry}}.
\label{eq:structure}
\end{equation}
Left side: the entropy-area density scaled by the Euclidean horizon
period, both set by flat spacetime geometry. Right side: the
gravitational coupling set by the Newtonian limit and the trace
structure of the Einstein tensor. The equation is satisfied if and
only if $\eta = 1/(4\ell_P^2)$. No observational input required.

Jacobson's machine --- \textit{specify an entropy functional, derive
a gravity theory} --- now fully specifies its own entropy functional.
The machine is self-contained.

\subsection{Why Area Scaling Is Geometrically Necessary}

The $S \propto A$ relation is not merely consistent with the causal
geometry --- it is the unique scaling required by it. Information from
bulk decoherence events can only be permanently encoded on the causal
boundary: anything inside the horizon is causally disconnected from
external observers. Bulk volume encoding would place information
inaccessible to those observers. Area scaling is therefore a causal
necessity~\citep{tHooft1993,Susskind1995}, and the coefficient
$1/(4\ell_P^2)$ is the unique value consistent with the Rindler and
Einstein geometric factors derived above.

\subsection{The Same Projection Structure at Cosmological Scale}

The ratio $2\pi/8\pi = 1/4$ involves the same geometric projection
factors --- a full-sphere solid angle ($4\pi$), its doubling ($8\pi$),
and a topological period ($2\pi$) --- that appear in the IAM derivation
of the cosmological constant~\citep{IAM_archive}. In that context, the
de~Sitter causal patch factor $2/\pi$ and its inverse $\pi/2$ (the
reverse holographic projection) appear independently in the CC
derivation and in the empirical Koide formula for lepton
masses~\citep{IAM_archive,Koide1983}, with their product being exactly
1. The present result identifies the Planck-scale analog: the ratio
$2\pi/8\pi$ at the horizon scale corresponds to the same holographic
projection structure operating at the fundamental causal boundary of
flat spacetime.

% ─────────────────────────────────────────────────────────────────────────────
\section{The Remaining Open Problem}
\label{sec:open}
% ─────────────────────────────────────────────────────────────────────────────

The present results close three questions that Jacobson left open:

\textbf{(A) Why does $S$ scale with $A$?} Derived in
Section~\ref{sec:area}: decoherence produces classical information
irreversibly; permanent information must be encoded at the causally
accessible 2D boundary; volume encoding is causally forbidden.
Area scaling is a necessary consequence of the causal structure.

\textbf{(B) Why is the coefficient $1/4$?} Derived in
Section~\ref{sec:factors}: $1/4 = 2\pi/8\pi$, where $2\pi$ is the
Euclidean Rindler cone angle and $8\pi$ is the Einstein normalization.
Both are geometric facts about flat spacetime. Neither assumes $S
\propto A$.

\textbf{(C1) Why does one bit occupy $4\ell_P^2$?} Derived in
Section~\ref{sec:connections}: the minimum horizon area disturbed
by one Landauer cost event at Planck-scale surface gravity is
$4\ell_P^2$, from the same $8\pi/2\pi = 4$ ratio. The entropy
coefficient and the minimum encoding area are the same geometric
fact.

The remaining open question is distinct from all three above. It is:

\begin{center}
\textit{Why does a minimum resolvable length $\ell_P = \sqrt{\hbar G/c^3}$
exist at all?}
\end{center}

$\ell_P^2 = \hbar G/c^3$ is the unique area constructible from the
three constants governing quantum mechanics ($\hbar$), gravity ($G$),
and relativity ($c$). If a minimum encoding area exists, it must be
proportional to $\ell_P^2$ by dimensional necessity --- no other
combination is available. But why a minimum area exists in the first
place, and why the three constants $\hbar$, $G$, $c$ combine to set
it, is a question about the discrete microstructure of spacetime at
the Planck scale. This lies beyond the thermodynamic level of
Jacobson's derivation and cannot be addressed from within it.

Loop quantum gravity and string theory approach this question from
different directions. The present derivation is neutral between those
approaches. It identifies $\ell_P^2$ as the minimum unit by dimensional
necessity, and derives all three geometric factors ($S \propto A$,
$\eta = 1/4$, one bit per $4\ell_P^2$) from causal spacetime
structure. The single remaining gap is the quantum gravity account of
why Planck-scale discreteness exists.


% ─────────────────────────────────────────────────────────────────────────────
\section{Discussion}
\label{sec:discussion}
% ─────────────────────────────────────────────────────────────────────────────

The result is best stated precisely rather than broadly.

Jacobson's 1995 paper made two assumptions: $S \propto A$ and
$\eta$ is a free parameter. Both were reasonable at the time.
$S \propto A$ was motivated by the Bekenstein--Hawking result for
black holes; $\eta$ was recovered by matching to Newton's constant.
What was not recognized is that both assumptions follow from the
causal geometry of flat spacetime.

$S \propto A$ follows from the irreversibility of decoherence and
the causal inaccessibility of the bulk: information produced by an
irreversible quantum-to-classical transition must be encoded at the
2D causal boundary because nowhere else is causally accessible to
external observers. This does not require quantum gravity; it
requires only the causal structure of spacetime and the
irreversibility of decoherence established by Landauer's
principle~\citep{Landauer1961} and Zurek's decoherence
framework~\citep{Zurek2003}.

$\eta = 1/(4\ell_P^2)$ follows from the ratio $2\pi/8\pi$ where
both factors are present in Jacobson's derivation: the $2\pi$ from
the Euclidean Rindler cone angle (Bisognano--Wichmann~\citep{Bisognano1976}),
and the $8\pi$ from the Einstein normalization (Wald~\citep{Wald1984}).
Jacobson had both. He did not divide them.

The minimum encoding area $4\ell_P^2 = \delta A_\mathrm{min}$
follows from the Landauer cost of one decoherence event at
Planck-scale surface gravity, with the same $8\pi/2\pi = 4$
geometric ratio. The product $\eta \cdot \delta A_\mathrm{min} = 1$
is exact --- one bit per minimum encoding area --- and is not a
separate result but the same result stated as a consistency check.

Three things Jacobson assumed are now derived. One thing remains
open: why Planck-scale discreteness exists at all. That is the
genuine quantum gravity question, and it is a much narrower and
more precisely stated gap than any existing treatment of the
Bekenstein--Hawking coefficient provides.

The thermodynamic level of the derivation of $S = A/(4\ell_P^2)$
is complete.

% ─────────────────────────────────────────────────────────────────────────────
\section*{Acknowledgments}
% ─────────────────────────────────────────────────────────────────────────────

The thermodynamic framework of T.\ Jacobson~\citep{Jacobson1995} is
the foundation of this work. This paper was motivated directly by a
personal response from Professor P.J.E.\ Peebles, who wrote:
\textit{``It is good to think independently, as you have done by
exploring consequences of the entanglement of the whole world, assuming
quantum physics is accurate enough to be trusted for this purpose.''}
That statement --- identifying the quantum-to-classical derivation as
the foundational gap requiring closure before the cosmological
framework could be trusted --- directly motivated the analysis
presented here. The Unruh--Rindler geometry of W.G.\ Unruh~\citep{Unruh1976}
and the parallel result of P.C.W.\ Davies~\citep{Davies1975} provide
the $2\pi$ factor. The Bisognano--Wichmann theorem~\citep{Bisognano1976}
establishes it in the full QFT setting.

% ─────────────────────────────────────────────────────────────────────────────
\section*{Data Availability}
% ─────────────────────────────────────────────────────────────────────────────

All derivations, numerical verifications, and companion papers are
publicly archived at \url{https://github.com/hmahaffeyges/IAM-Validation}
with canonical Zenodo DOI
\href{https://doi.org/10.5281/zenodo.18702042}{10.5281/zenodo.18702042}.

% ─────────────────────────────────────────────────────────────────────────────

\begin{thebibliography}{99}

\bibitem[Zurek(2003)]{Zurek2003}
Zurek, W.H. (2003).
Decoherence, einselection, and the quantum origins of the classical.
\textit{Rev.\ Mod.\ Phys.}, 75, 715--775.
\url{https://doi.org/10.1103/RevModPhys.75.715}

\bibitem[Bekenstein(1973)]{Bekenstein1973}
Bekenstein, J.D. (1973).
Black holes and entropy.
\textit{Phys.\ Rev.\ D}, 7, 2333--2346.
\url{https://doi.org/10.1103/PhysRevD.7.2333}

\bibitem[Bisognano \& Wichmann(1976)]{Bisognano1976}
Bisognano, J.J. and Wichmann, E.H. (1976).
On the duality condition for quantum fields.
\textit{J.\ Math.\ Phys.}, 17, 303--321.
\url{https://doi.org/10.1063/1.522898}

\bibitem[Cai \& Kim(2005)]{CaiKim2005}
Cai, R.-G. and Kim, S.P. (2005).
First law of thermodynamics and Friedmann equations of
Friedmann--Robertson--Walker universe.
\textit{J.\ High Energy Phys.}, 02, 050.
\url{https://doi.org/10.1088/1126-6708/2005/02/050}

\bibitem[Davies(1975)]{Davies1975}
Davies, P.C.W. (1975).
Scalar particle production in Schwarzschild and Rindler metrics.
\textit{J.\ Phys.\ A}, 8, 609--616.
\url{https://doi.org/10.1088/0305-4470/8/4/022}

\bibitem[Gibbons \& Hawking(1977)]{Gibbons1977}
Gibbons, G.W. and Hawking, S.W. (1977).
Cosmological event horizons, thermodynamics, and particle creation.
\textit{Phys.\ Rev.\ D}, 15, 2738--2751.
\url{https://doi.org/10.1103/PhysRevD.15.2738}

\bibitem[Hawking(1975)]{Hawking1975}
Hawking, S.W. (1975).
Particle creation by black holes.
\textit{Commun.\ Math.\ Phys.}, 43, 199--220.
\url{https://doi.org/10.1007/BF02345020}

\bibitem[IAM archive(2026)]{IAM_archive}
Mahaffey, H.W. (2026).
IAM Validation Repository.
Zenodo. \url{https://doi.org/10.5281/zenodo.18702042}.
\url{https://github.com/hmahaffeyges/IAM-Validation}

\bibitem[Jacobson(1995)]{Jacobson1995}
Jacobson, T. (1995).
Thermodynamics of spacetime: the Einstein equation of state.
\textit{Phys.\ Rev.\ Lett.}, 75, 1260--1263.
\url{https://doi.org/10.1103/PhysRevLett.75.1260}

\bibitem[Koide(1983)]{Koide1983}
Koide, Y. (1983).
New view of quark and lepton mass hierarchy.
\textit{Phys.\ Rev.\ D}, 28, 252--254.
\url{https://doi.org/10.1103/PhysRevD.28.252}

\bibitem[Landauer(1961)]{Landauer1961}
Landauer, R. (1961).
Irreversibility and heat generation in the computing process.
\textit{IBM J.\ Res.\ Dev.}, 5, 183--191.
\url{https://doi.org/10.1147/rd.53.0183}

\bibitem[Padmanabhan(2010)]{Padmanabhan2010}
Padmanabhan, T. (2010).
Thermodynamical aspects of gravity: new insights.
\textit{Rep.\ Prog.\ Phys.}, 73, 046901.
\url{https://doi.org/10.1088/0034-4885/73/4/046901}

\bibitem[Rovelli \& Smolin(1995)]{Rovelli1995}
Rovelli, C. and Smolin, L. (1995).
Discreteness of area and volume in quantum gravity.
\textit{Nucl.\ Phys.\ B}, 442, 593--622.
\url{https://doi.org/10.1016/0550-3213(95)00150-Q}

\bibitem[Strominger \& Vafa(1996)]{Strominger1996}
Strominger, A. and Vafa, C. (1996).
Microscopic origin of the Bekenstein--Hawking entropy.
\textit{Phys.\ Lett.\ B}, 379, 99--104.
\url{https://doi.org/10.1016/0370-2693(96)00345-0}

\bibitem[Susskind(1995)]{Susskind1995}
Susskind, L. (1995).
The world as a hologram.
\textit{J.\ Math.\ Phys.}, 36, 6377--6396.
\url{https://doi.org/10.1063/1.531249}

\bibitem['t~Hooft(1993)]{tHooft1993}
't~Hooft, G. (1993).
Dimensional reduction in quantum gravity.
\textit{arXiv:gr-qc/9310026}.

\bibitem[Unruh(1976)]{Unruh1976}
Unruh, W.G. (1976).
Notes on black-hole evaporation.
\textit{Phys.\ Rev.\ D}, 14, 870--892.
\url{https://doi.org/10.1103/PhysRevD.14.870}

\bibitem[Wald(1984)]{Wald1984}
Wald, R.M. (1984).
\textit{General Relativity}.
University of Chicago Press.

\end{thebibliography}

\end{document}
